\scanpage{77}
\mspage{34}
\uline{15}. Introduction of the prospaces :
${\unsure{V|\mathcal V}}_{i_0.\!-.i_n},\ {\unsure{V|\mathcal V}}_J^{\deleted{\unsure{*}}L}$,
% The title has a dot, a short horizontal dash, and a dot between the indices.
% Both V forms have a curved bowl; plain/calligraphic alphabet assignment remains open.

We denote by ${\unsure{V|\mathcal V}}_i^{\unsure{*|<|\deleted{*}}}$
\unsure{\deleted{the}|the} tubular \inserted{pro}neighborhood
(in the sense of prospaces) of $X_i^*=X_i\setminus\bigcup_{j<i}X_j$
in $X$. Thus, \deleted{\unsure{we}} \unsure{we have}
% The raised pro has a curved leader ending immediately before voisinage.
\[
(51)\qquad X_i^*\longrightarrow {\unsure{V|\mathcal V}}_i^{\unsure{*|\deleted{*}}}
\qquad\begin{gathered}
\text{homotopy equiva-}\\
\text{lence \unsure{of|of} prospaces.}
\end{gathered}
\]
% The hyphen after homo is written at the right paper edge.

\deleted{\unsure{For a flag}}
If\unsure{)} $i<j$, we set
% The separate curved stroke after Si is retained; no opening parenthesis is supplied.
\[
(52)\qquad
\unsure{\rlap{$X$}V|\rlap{$X$}\mathcal V}_{ij}={\unsure{V|\mathcal V}}_i\cap {\unsure{V|\mathcal V}}_j
\]
% The left base in (52) is one superposed X/V cluster with a single shared ij.
% The overprinting order and deletion state cannot be separated; no second index is supplied.

\margin{For a
\diagram{\begin{tikzcd}[column sep=small]
{\unsure{V|\mathcal V}}_{ij}^{\unsure{b}}={\unsure{V|\mathcal V}}_i\cap X_j^*\arrow[r,hook]&{}
\end{tikzcd}}
is this a homotopy equivalence ?}
% The oblique left note ends its formula at the curved arrowhead; no target letters are visible.

\margin{NB ${\unsure{V|\mathcal V}}_{ij}$ makes sense \deleted{\unsure{if}} if $i\neq j$
but if $i$ and $j$ are not comparable, we have
${\unsure{V|\mathcal V}}_i\cap {\unsure{V|\mathcal V}}_j=\varnothing$,
\inserted{(at least if the open subsets of $X$ are normal, for example $X$ metrizable)}
if $i=j$, we have ${\unsure{V|\mathcal V}}_{ii}={\unsure{V|\mathcal V}}_i$}
% The long curved leader joins the parenthetical insertion to the intersection-empty clause.
% The parenthetical words are physically below V_ii=V_i; their attachment precedes si i=j.

so that we have
\diagram{\begin{tikzcd}[column sep=large,row sep=huge]
&{\unsure{V|\mathcal V}}_{ij}\arrow[r,hook,"b_{ij}"]\arrow[d,hook,"s_{ij}"]&{\unsure{V|\mathcal V}}_j\\
\Bigl(X_i
\arrow[r,shift right=1.2ex]
\arrow[r,phantom,"{\substack{\unsure{(w)|(h)}\\\sim\\
\unsure{\longleftarrow|\deleted{\longleftarrow}}}}"]
&\Bigr){\unsure{V|\mathcal V}}_i&
\end{tikzcd}}
% Four written connectors: two hooked main arrows, the dense left-pointing lower mark,
% and a separate lower right-pointing arrow. The left-pointing mark has unresolved strike state.
% The closing parenthesis lies before V_i, with the one tilde and (w)/(h) above the lower marks.
We have :
\[
(53)\qquad X=\varinjlim\bigl({\unsure{V|\mathcal V}}_i,{\unsure{V|\mathcal V}}_{ij},
\unsure{\deleted{\}}|\}}\ b_{ij},s_{ij}\bigr)
\]
% The brace-like glyph between the comma and b is crossed/reworked; its strike state remains open.
$X$ is obtained by \uline{gluing} the ${\unsure{V|\mathcal V}}_i$
along the ${\unsure{V|\mathcal V}}_{ij}$.

We also define, more generally, the
\[
(54)\qquad {\unsure{V|\mathcal V}}_{i_0.\!-i_n}={\unsure{V|\mathcal V}}_{i_0}\cap
\mathord{-}\,\mathord{-}\cap {\unsure{V|\mathcal V}}_{i_n}
\]
% The left index has dot-dash; the right continuation has two separate dashes.
% The final V is retraced, without a recoverable independently cancelled underlying letter.
for a flag $i_0<i_1<.\!-<i_n$
\margin{\uline{NB} The expression makes sense $=$ if
the set $\{i_0,\mathord{-}i_n\}$ does not form a flag,
but it is empty if it is not a flag}
% The bottom marginal note is oblique. Its = has two horizontal strokes, the lower slightly wavy.
% The continuation within the written braces is one dash, with no second comma before i_n.
% Pale bottom-left foliation in another, unnamed hand:
\margin{76}
